\documentclass[12pt]{article}
\usepackage[utf8]{inputenc}
\usepackage{amsmath}
\usepackage{amssymb}
\usepackage{amsfonts}
\usepackage{geometry}
\usepackage{enumitem}
\usepackage{chemfig}
\usepackage{multicol}
\geometry{a4paper, margin=1in}

\begin{document}

\section*{Exam Review Questions}

The following is a small selection of questions that can help you prepare for the final exam. This is \underline{not} a complete exam review sheet. You will be required to consult various reference sheets for data to help solve a few of the problems.

\subsection*{Structure \& Properties Unit}

\begin{enumerate}[label=\arabic*.]
    \item What is the letter code for a subshell with $\ell = 3$?
    \vspace{1.5cm}

    \item How many orbitals are found in the $f$ subshell?
    \vspace{1.5cm}

    \item The ground state electron configuration of a phosphorus atom is $\text{1s}^2\,\text{2s}^2\,\text{2p}^6\,\text{3s}^2\,\text{3p}^3$. More specifically, $\text{3p}^3$ can be expanded to: $\text{3p}_x\,\text{3p}_y\,\text{3p}_z$. What does this mean?
    \vspace{3cm}

    \item What feature of their ground state electron configurations distinguishes the transition elements from other elements?
    \vspace{3cm}

    \item Explain the term isoelectronic. Provide an example.
    \vspace{3cm}

    \item Explain the term polar bond.
    \vspace{2.5cm}

    \item What is the maximum number of electrons allowed within the $n = 4$ region of space? What is the maximum number of electrons belonging to the principal quantum number $n = 4$?
    \vspace{3cm}

    \item Write the longhand ground state electron configuration for boron \& write the shorthand ground state electron configuration for sulphur.
    \vspace{3cm}

    \item Draw energy level diagrams to show the ground state electron configurations of argon \& zinc.
    \vspace{5cm}

    \item How many electrons are in the valence energy level of bromine \& xenon?
    \vspace{2.5cm}

    \item Name the neutral element that has the following ground state electron configuration: $\text{1s}^2\,\text{2s}^2\,\text{2p}^6\,\text{3s}^2\,\text{3p}^6\,\text{4s}^2\,\text{3d}^2$.
    \vspace{2cm}

    \item What is a resonance structure? Draw all resonance structures for the carbonate ion, $\text{CO}_3^{2-}$.
    \vspace{4.5cm}

    \item Explain why $\text{NH}_3$ is trigonal pyramidal whereas $\text{BCl}_3$ is trigonal planar.
    \vspace{3.5cm}

    \item Using VSEPR theory, explain why the bond angle in $\text{CO}_2$ is $180^\circ$.
    \vspace{3cm}

    \item Draw the Lewis structures for $\text{AsCl}_5$ \& $\text{BrO}_3^-$. Name the shape of these structures. Indicate if the molecules are polar \& the type of intermolecular forces that would result between 2 molecules of $\text{AsCl}_5$ \& between 2 molecules of $\text{BrO}_3^-$.
    \vspace{5cm}

    \item Explain why gases may be compressed while solids \& liquids may not be compressed to any extent.
    \vspace{3cm}

    \item What is coordinate covalent bonding?
    \vspace{2.5cm}

    \item For the element titanium, $\text{Ti}$:
    \begin{enumerate}[label=\alph*)]
        \item Predict the most stable electron configuration for neutral titanium.
        \vspace{2cm}
        \item Without looking at the periodic table for the actual charges, predict the 2 most stable ionic charges for titanium. Which one should be more stable?
        \vspace{2.5cm}
        \item On the periodic table, it says $\text{Ti}^{3+}$ is stable, but is not the most stable ionic form. Provide a possible electron configuration that would explain why $\text{Ti}^{3+}$ is considered a stable ionic charge.
        \vspace{3cm}
    \end{enumerate}

    \item Draw \& explain the valence bonding \& hybridization that occurs in ethane, ethene \& ethyne.
    \vspace{5cm}

    \item Which will have the higher boiling point?
    \begin{enumerate}[label=\alph*)]
        \item $\text{SF}_2$ or $\text{SBr}_2$
        \vspace{1.5cm}
        \item $\text{Cl}_2$ or $\text{I}_2$
        \vspace{1.5cm}
        \item $\text{Cu}$ or $\text{C}_3\text{H}_8$
        \vspace{1.5cm}
    \end{enumerate}

    \item What kind of crystal solid is formed from $\text{MgO}$?
    \vspace{2cm}
\end{enumerate}

\newpage

\subsection*{Energy Changes \& Rates of Reactions Unit}

\begin{enumerate}[label=\arabic*.]
    \item How much energy is released when $15.0\text{ grams}$ of water vapour at $112.0^\circ\text{C}$ becomes $-15.0^\circ\text{C}$ solid ice?
    \vspace{4.5cm}

    \item Calculate $\Delta H^\circ$ for the following reaction: 
    \[
    2\,\text{HCl(g)} + \text{CaO(s)} \rightarrow \text{CaCl}_2\text{(s)} + \text{H}_2\text{O(g)}
    \]
    using enthalpies of formation values.
    \vspace{4.5cm}

    \item For the reaction: $3\,\text{H}_2 + \text{N}_2 \rightleftharpoons 2\,\text{NH}_3$
    \begin{enumerate}[label=\alph*)]
        \item How does the rate of disappearance of hydrogen compare to the rate of disappearance of nitrogen?
        \vspace{2cm}
        \item How does the rate of appearance of ammonia compare to the rate of disappearance of nitrogen?
        \vspace{2cm}
        \item If the rate of the appearance of $\text{NH}_3$ is $2.5 \times 10^{-6}\text{ mol/(L}\cdot\text{s)}$, how fast are $\text{H}_2$ \& $\text{N}_2$ being used?
        \vspace{3cm}
    \end{enumerate}

    \item The rate law for the reaction $2\,\text{NO} + \text{O}_2 \rightarrow 2\,\text{NO}_2$ is $r = k[\text{NO}]^2[\text{O}_2]$. At $25^\circ\text{C}$, $k = 7.1 \times 10^9\text{ L}^2/(\text{mol}^2\cdot\text{s})$.
    \begin{enumerate}[label=\alph*)]
        \item What is the rate of reaction when $[\text{NO}] = 0.0010\text{ mol/L}$ \& $[\text{O}_2] = 0.034\text{ mol/L}$?
        \vspace{2.5cm}
        \item Is the reaction $2\,\text{NO} + \text{O}_2 \rightarrow 2\,\text{NO}_2$ an elementary step? Explain why.
        \vspace{2.5cm}
    \end{enumerate}

    \item Write the equation for the formation of the explosive TNT (trinitrotoluene), $\text{C}_7\text{H}_5\text{O}_6\text{N}_3\text{(s)}$, from its elements at standard state.
    \vspace{3.5cm}

    \item When $5.60\text{ kJ}$ of thermal energy is added to $167\text{ grams}$ of gaseous ammonia, $\text{NH}_3$, at $45.0^\circ\text{C}$, the temperature of the gas rises to $60.0^\circ\text{C}$. From this data determine the following:
    \begin{enumerate}[label=\alph*)]
        \item the specific heat capacity of $\text{NH}_3$
        \vspace{2.5cm}
        \item the molar heat capacity of $\text{NH}_3$
        \vspace{2.5cm}
    \end{enumerate}

    \item From the following three thermochemical equations, calculate the enthalpy change for the reaction: $\text{FeO} + \text{CO} \rightarrow \text{Fe} + \text{CO}_2$
    \begin{align*}
    \text{Fe}_2\text{O}_3 + 3\,\text{CO} &\rightarrow 2\,\text{Fe} + 3\,\text{CO}_2 & \Delta H^\circ &= -25\text{ kJ} \\
    3\,\text{Fe}_2\text{O}_3 + \text{CO} &\rightarrow 2\,\text{Fe}_3\text{O}_4 + \text{CO}_2 & \Delta H^\circ &= -47\text{ kJ} \\
    \text{Fe}_3\text{O}_4 + \text{CO} &\rightarrow 3\,\text{FeO} + \text{CO}_2 & \Delta H^\circ &= +38\text{ kJ}
    \end{align*}
    \vspace{5cm}

    \item A copper container has a mass of $305\text{ grams}$ \& contains $255\text{ grams}$ of water. When $1.01\text{ grams}$ of propanol, $\text{C}_3\text{H}_8\text{O}$, is burned in the container, the copper \& contents increase in temperature by $28.6^\circ\text{C}$. Calculate the molar enthalpy of combustion of propanol using this experimental data.
    \vspace{4.5cm}

    \item A $50.0\text{ gram}$ piece of aluminum is heated to $100^\circ\text{C}$ \& then put into a beaker containing $150\text{ mL}$ of water at $20.0^\circ\text{C}$. Assuming no loss of thermal energy to the surroundings, calculate the final temperature of the water. Describe what is happening on the molecular level.
    \vspace{4.5cm}

    \item A $4.05\text{ gram}$ sample of sulphur is burned in excess oxygen. The energy produced raises the temperature of $1.00\text{ L}$ of water from $23.15^\circ\text{C}$ to $32.03^\circ\text{C}$. From this information calculate $\Delta H^\circ_f$ of sulphur dioxide.
    \vspace{4.5cm}

    \item The initial rate of production of bromine was measured using different initial concentrations of each reactant at $25.0^\circ\text{C}$. The equation \& results are summarized below:
    \[
    4\,\text{HBr(g)} + \text{O}_2\text{(g)} \rightarrow 2\,\text{Br}_2\text{(l)} + 2\,\text{H}_2\text{O(g)}
    \]
    \begin{center}
    \begin{tabular}{|c|c|c|}
    \hline
    \textbf{Initial Rate of Reaction (mol/(L$\cdot$s))} & \textbf{Initial [HBr] (mol/L)} & \textbf{Initial [$\text{O}_2$] (mol/L)} \\ \hline
    0.0042 & 0.010 & 0.010 \\ \hline
    0.0084 & 0.010 & 0.020 \\ \hline
    0.0126 & 0.030 & 0.010 \\ \hline
    \end{tabular}
    \end{center}
    \begin{enumerate}[label=\alph*)]
        \item Determine the relationship between rate of decomposition \& the concentration of each reactant.
        \vspace{2cm}
        \item Determine the rate law for the reaction \& the overall order of the reaction.
        \vspace{2cm}
        \item Determine the value of the rate constant, $k$. Include units.
        \vspace{2cm}
        \item Calculate $\Delta H^\circ$ for the reaction as written using $\Delta H^\circ_f$ values. Is the reaction endothermic or exothermic?
        \vspace{2.5cm}
        \item Draw the potential energy diagram for the reaction. Label all critical points \& values.
        \vspace{4.5cm}
        \item What would happen if the reaction occurred at $30.0^\circ\text{C}$? Explain why.
        \vspace{2.5cm}
    \end{enumerate}
\end{enumerate}

\newpage

\subsection*{Chemical Systems \& Equilibrium Unit}

\begin{enumerate}[label=\arabic*.]
    \item For the following reaction: $\text{Br}_2\text{(l)} + 5\,\text{F}_2\text{(g)} \rightleftharpoons 2\,\text{BrF}_5\text{(g)}$, $K_{eq} = 1.0 \times 10^{18}$.
    \begin{enumerate}[label=\alph*)]
        \item Write the equilibrium law expression for the reaction.
        \vspace{1.5cm}
        \item At equilibrium, is the reaction mostly reactants or products dominant?
        \vspace{1.5cm}
        \item What should you do to the equation if you want to apply the ``100'' rule assumption?
        \vspace{1.5cm}
        \item If $0.12\text{ mol/L}$ of $\text{BrF}_5$ is initially present, calculate the equilibrium concentrations of all compounds.
        \vspace{3.5cm}
        \item If the reaction is: $2\,\text{Br}_2\text{(l)} + 10\,\text{F}_2\text{(g)} \rightleftharpoons 4\,\text{BrF}_5\text{(g)}$, what is the value of $K_{eq}$?
        \vspace{2cm}
    \end{enumerate}

    \item At $1000^\circ\text{C}$ methane reacts with water as follows: $\text{CH}_4 + \text{H}_2\text{O} \rightleftharpoons \text{CO} + 3\,\text{H}_2$, all compounds are gaseous. In one experiment, the equilibrium concentrations of the gases were $[\text{CH}_4] = 2.97 \times 10^{-3}\text{ mol/L}$, $[\text{H}_2\text{O}] = 7.94 \times 10^{-3}\text{ mol/L}$, $[\text{CO}] = 5.45 \times 10^{-3}\text{ mol/L}$, \& $[\text{H}_2] = 2.1 \times 10^{-3}\text{ mol/L}$. Calculate $K_{eq}$ at this temperature.
    \vspace{3.5cm}

    \item For the following process: $2\,\text{FeBr}_3\text{(s)} \rightleftharpoons 2\,\text{FeBr}_2\text{(s)} + \text{Br}_2\text{(g)}$, the equilibrium constant, $K_{eq}$, is $0.98$. A $3.0\text{ L}$ reaction vessel contains $1.2\text{ mol}$ of $\text{FeBr}_3$, $3.6\text{ mol}$ of $\text{FeBr}_2$ \& $2.1\text{ mol}$ of $\text{Br}_2$ at this temperature. Is the system at equilibrium?
    \vspace{3.5cm}

    \item The equilibrium constant for the gaseous system $\text{SO}_3 + \text{NO} \rightleftharpoons \text{NO}_2 + \text{SO}_2$ was found to be $0.250$ at a certain temperature. If $0.300\text{ mol}$ of $\text{SO}_3$ \& $0.300\text{ mol}$ of $\text{NO}$ were placed in a $2.00\text{ L}$ container \& allowed to react, what would be the equilibrium concentrations of each gas?
    \vspace{4.5cm}

    \item The gaseous system $2\,\text{HCl} + \text{I}_2 \rightleftharpoons 2\,\text{HI} + \text{Cl}_2$ has $K_{eq} = 1.6 \times 10^{-34}$ at $25^\circ\text{C}$. Suppose $1.00\text{ mol}$ of $\text{HCl}$ \& $1.00\text{ mol}$ of $\text{I}_2$ was placed in a $1.00\text{ L}$ container. What would be the equilibrium concentrations of $\text{HI}$ \& $\text{Cl}_2$ in the container?
    \vspace{4.5cm}

    \item A salt whose formula is of the form $\text{MX}$ has $K_{sp} = 2.0 \times 10^{-10}$. Another slightly soluble salt, $\text{MX}_3$ must have what $K_{sp}$ value if the molar solubilities of the two salts are identical?
    \vspace{4cm}

    \item Draw the concentration vs time graph illustrating the following situations for $\text{N}_2\text{(g)} + \text{O}_2\text{(g)} \rightleftharpoons 2\,\text{NO(g)}$; $\Delta H = +179\text{ kJ}$
    \begin{enumerate}[label=\alph*)]
        \item initially start with $[\text{NO}]$ only
        \vspace{1.5cm}
        \item increase the partial pressure of nitrogen gas at constant volume
        \vspace{1.5cm}
        \item decrease the volume of the container
        \vspace{1.5cm}
        \item add $\text{He(g)}$ at constant pressure
        \vspace{1.5cm}
        \item increase the temperature of the system
        \vspace{1.5cm}
    \end{enumerate}

    \item What are the ion concentrations in a saturated solution of silver sulphate, $\text{Ag}_2\text{SO}_4$?
    \vspace{3cm}

    \item For the insoluble calcium fluoride, $\text{CaF}_2$, if the solubility is $0.125\text{ g}/100\text{ mL}$ at $20^\circ\text{C}$:
    \begin{enumerate}[label=\alph*)]
        \item Calculate the $K_{sp}$ value for $\text{CaF}_2$ at $20^\circ\text{C}$.
        \vspace{3cm}
        \item Will a precipitate of $\text{CaF}_2$ form when $0.084\text{ grams}$ of sodium fluoride, $\text{NaF}$, is dissolved in $1.00\text{ L}$ of a $0.010\text{ mol/L}$ aqueous solution of calcium chloride, $\text{CaCl}_2$?
        \vspace{3.5cm}
    \end{enumerate}

    \item Calculate the maximum iodide ion concentration that can be added to $0.010\text{ mol/L}$ in lead (II) nitrate, $\text{Pb(NO}_3)_2$ before precipitation occurs.
    \vspace{3.5cm}

    \item A $1.50\text{ L}$ solution of $0.250\text{ mol/L}$ aqueous sodium hydroxide, $\text{NaOH}$, is added to $1.00\text{ L}$ solution of $0.1875\text{ mol/L}$ magnesium nitrate, $\text{Mg(NO}_3)_2$. Magnesium hydroxide, $\text{Mg(OH)}_2$, precipitates from the mixture. Calculate the:
    \begin{enumerate}[label=\alph*)]
        \item maximum amount of $\text{Mg(OH)}_2$ (in mol) that can form (assuming it cannot dissolve at all).
        \vspace{3cm}
        \item mass of magnesium hydroxide that precipitates.
        \vspace{2.5cm}
        \item pH for the solution.
        \vspace{3.5cm}
    \end{enumerate}

    \item Aspirin, $\text{HC}_9\text{H}_7\text{O}_4$, is a monoprotic acid whose $K_a = 3.27 \times 10^{-4}$. Does a solution of the copper (II) salt of aspirin in water test acidic, basic or neutral? In other words, is $\text{Cu(C}_9\text{H}_7\text{O}_4)_2$ acidic, basic or neutral?
    \vspace{3.5cm}

    \item Let's assume ascorbic acid is monoprotic. If $1.76\text{ grams}$ of ascorbic acid, $\text{C}_6\text{H}_8\text{O}_6$, was placed in $100.0\text{ mL}$ water. What is the pH of the solution?
    \vspace{3.5cm}

    \item An aqueous solution containing $0.302\text{ grams}$ of a triprotic acid, $\text{H}_3\text{X}$, requires $39.30\text{ mL}$ of a $0.120\text{ mol/L}$ solution of sodium hydroxide, $\text{NaOH}$, for the neutralization of all three protons. Calculate the molecular mass of the acid.
    \vspace{3.5cm}

    \item In an experiment, $0.030\text{ L}$ of $0.25\text{ mol/L}$ ethanoic acid $\text{CH}_3\text{COOH}$, was added to $0.020\text{ L}$ of $0.23\text{ mol/L}$ sodium hydroxide, $\text{NaOH}$. Find the pH of the resultant solution. \textit{Hint: the base will neutralize only part of the acid. What is the pH of the leftover acid?}
    \vspace{4.5cm}

    \item What is the pH of a buffer solution that contains $0.10\text{ mol/L}$ $\text{NaHCO}_3$ \& $0.12\text{ mol/L}$ $\text{H}_2\text{CO}_3$?
    \vspace{3.5cm}
\end{enumerate}

\newpage

\subsection*{Electrochemistry Unit}

\begin{enumerate}[label=\arabic*.]
    \item What is the difference between oxidation numbers \& valence charges?
    \vspace{2.5cm}

    \item What are the oxidation numbers in the following: sulphur in $\text{Cu}_2\text{SO}_4$ \& carbon in $\text{C}_2\text{H}_3\text{O}_2$?
    \vspace{2.5cm}

    \item Identify the reactant that is oxidized \& the reactant that is reduced in each reaction. Which reactant is the oxidizing agent?
    \[
    \text{H}_2\text{C}_2\text{O}_4 + \text{H}_2\text{O}_2 \rightarrow 2\,\text{CO}_2 + 2\,\text{H}_2\text{O}
    \]
    \vspace{3cm}

    \item Balance the following redox reactions:
    \begin{enumerate}[label=\alph*)]
        \item $\text{Ag} + \text{Cr}_2\text{O}_7^{2-} \rightarrow \text{Cr}^{3+} + \text{Ag}^+$ \quad (acidic solution)
        \vspace{3cm}
        \item $\text{Al} + \text{MnO}_4^- \rightarrow \text{MnO}_2 + \text{Al}^{3+}$ \quad (basic solution)
        \vspace{3cm}
    \end{enumerate}

    \item How can a solution of aqueous metal cations be isolated \& purified in a metal plating process?
    \vspace{3.5cm}

    \item Which of the following is the best oxidizing agent in aqueous solution: $\text{Cl}_2$, $\text{Cu}^{2+}$ or $\text{Ag}^+$?
    \vspace{2.5cm}

    \item Determine if the following unbalanced reaction is spontaneous or not: $\text{Cd}^{2+} + \text{Fe} \rightarrow \text{Cd} + \text{Fe}^{2+}$.
    \vspace{2.5cm}

    \item Calculate the minimum cell voltage of a lithium iodide, $\text{LiI}$, electrolytic cell, where iodine must oxidize.
    \vspace{3.5cm}

    \item Draw \& fully label the galvanic cell with the following notation: $\text{Al(s)} \mid \text{Al}^{3+}\text{(aq)} \parallel \text{Cu}^{2+}\text{(aq)} \mid \text{Cu(s)}$. Include appropriate chemicals, as needed. Include the direction of electron \& ion flow. Calculate $\Delta E$.
    \vspace{5cm}

    \item An aqueous solution of silver nitrate, $\text{AgNO}_3$, is subjected to electrolysis for $30.0\text{ minutes}$ using a constant current of $1.50\text{ amperes}$. Calculate the mass of silver metal deposited at the cathode.
    \vspace{3.5cm}

    \item Calculate the standard cell potential ($\Delta E$) at $25^\circ\text{C}$ for the reaction: $\text{Fe}^{2+}\text{(aq)} + \text{Ag}^+\text{(aq)} \rightleftharpoons \text{Fe}^{3+}\text{(aq)} + \text{Ag(s)}$.
    \vspace{3.5cm}

    \item Aluminum metal is produced by the electrolysis of molten aluminum oxide, $\text{Al}_2\text{O}_3$. Consider one such cell with a potential difference between the anode \& cathode of $5.00\text{ V}$. Electrons are supplied at a rate of $1.00 \times 10^5\text{ amperes}$. How much time (in hours) will it take for one cell to produce $1100\text{ kg}$ of aluminum?
    \vspace{4.5cm}

    \item Write the spontaneous reaction that occurs when $\text{Ni(ClO}_3)_2$ is mixed with $\text{CrCl}_2$ in the presence of $\text{NaOH}$.
    \vspace{3.5cm}
\end{enumerate}

\newpage

\subsection*{Organic Chemistry Unit}

\begin{enumerate}[label=\arabic*.]
    \item Which of alkanes or alkynes (assuming both have the same number of carbon atoms) have the higher boiling point \& why?
    \vspace{3cm}

    \item Which of ethers or alcohols (assuming both have the same number of carbon atoms) have the higher solubility in oil \& why?
    \vspace{3cm}

    \item Draw all the different structural isomers of $\text{C}_5\text{H}_8$.
    \vspace{5cm}

    \item What is the difference between a structural \& a geometric isomer?
    \vspace{3cm}

    \item Synthesize the following compounds (using the shortest number of steps) from the indicated starting materials. Use appropriate reagents in your synthesis.
    \begin{enumerate}[label=\alph*)]
        \item \quad \chemfig{*6(-=-(-C(=[1]O)-[-1]OH)=-=)} \quad from $\text{CH}_4$ \& \quad \chemfig{*6(-=-=-=)}
        \vspace{3.5cm}
        \item \quad \chemfig{-C(=[1]O)-[-1]O-} \quad from $\text{CH}_4$ \& $\text{C}_2\text{H}_6$
        \vspace{3.5cm}
        \item \quad \chemfig{-[1](-[3])-NH-[-1]} \quad from $\text{CH}_4$ \& $\text{C}_2\text{H}_6$
        \vspace{3.5cm}
        \item \quad \chemfig{-[1]=_[-1]} \quad from \quad \chemfig{-[1]-[-1]}
        \vspace{3.5cm}
    \end{enumerate}

    \item Predict the products for the following reactions. If no reaction occurs, write NR.
    \begin{enumerate}[label=\alph*)]
        \item $\text{CH}_4$ is mixed with chlorine gas in the presence of UV light
        \vspace{2cm}
        \item $\text{CH}_3\text{OH}$ is mixed with \quad \chemfig{-C(=[1]O)-[-1]OH} \quad in the presence of sulphuric acid
        \vspace{2.5cm}
        \item \quad \chemfig{*6(-=-=-=)} \quad is mixed with $\text{CH}_4$ in the presence of sulphuric acid
        \vspace{2.5cm}
        \item \quad \chemfig{-[1]=_[-1]} \quad is mixed with hydrogen chloride
        \vspace{2.5cm}
        \item \quad \chemfig{O*3(--)} \quad is mixed with an acid solution
        \vspace{2.5cm}
    \end{enumerate}

    \item What is a polymer? Provide an example of a condensation \& addition polymer.
    \vspace{3.5cm}
\end{enumerate}

\end{document}